\chapter{Identitas Program Studi}

\begin{tabular}{L{5cm}L{9cm}}
Program Studi (PS) & : D4 Teknik Informatika \\
Jurusan/Departemen & : Teknologi Informasi \\
Fakultas & : - \\
Perguruan Tinggi & : Politeknik Negeri Malang \\
Nomor SK pendirian PS & : SK DIKTI No 50/D/O/2010 \\
Tanggal SK pendirian PS & : 21 Mei 2010 \\
Pejabat Penandatangan SK & : Menteri Pendidikan Nasional a.n. Direktur Jenderal Pendidikan Tinggi \\
Bulan \& Tahun Dimulainya Pendidikan & : Mei 2010 \\
Nomor SK Izin Operasional & : 50/D/O/2010 \\
Tanggal SK Izin Operasional & : 21 Mei 2010 \\
Nilai Akreditasi Terakhir & : B \\
Nomor SK BAN-PT/LAM & : 1810/SK/BAN-PT/AkRED/DiplIV/VII/2018 \\
Alamat PS & : Jalan Soekarno-Hatta No. 9 Malang 65141, Po Box 04 Malang \\
No. Telepon PS & : +62 (0341) 404424 - 404425 \\
No. Faksimili PS & : +62 (0341) 404420 \\
Homepage dan E-mail PS & : \url{http://jti.polinema.ac.id}, \url{ti@polinema.ac.id} \\
Jenjang Pendidikan & : Diploma IV \\
Gelar Lulusan & : S.Tr.Kom. (Sarjana Terapan Komputer) \\
\end{tabular}

\clearpage

\section{Visi, Misi, dan Tujuan Jurusan Teknologi Informasi}

\subsection{Visi}
Jurusan Teknologi Informasi POLINEMA sebagai pusat unggulan di bidang teknologi informasi dan rekayasa perangkat lunak di tingkat nasional maupun internasional.

\subsection{Misi}
\begin{enumerate}
\item Melaksanakan pendidikan vokasi yang inovatif berdasarkan pada sistem pendidikan terapan dengan memanfaatkan kemajuan teknologi informasi dan telekomunikasi, sehingga mampu menghasilkan lulusan yang siap kerja dengan daya saing global.
\item Melaksanakan penelitian terapan berbasis produk dan jasa bidang informatika.
\item Melaksanakan pengabdian masyarakat dengan menggunakan kemajuan teknologi informasi untuk meningkatkan kesejahteraan.
\item Mewujudkan kerja sama yang saling menguntungkan dengan berbagai pihak baik di dalam maupun di luar negeri pada bidang teknologi informasi.
\end{enumerate}

\subsection{Tujuan}
\begin{enumerate}
\item Menghasilkan lulusan bidang teknologi informasi dan rekayasa perangkat lunak yang berketuhanan, beretika dan bermoral baik, berpengetahuan dan berketrampilan tinggi, siap bekerja dan/atau berwirausaha yang mampu bersaing dalam skala global.
\item Menghasilkan penelitian terapan bidang teknologi informasi dan rekayasa perangkat lunak yang berskala internasional, meningkatkan efektifitas, efisiensi, dan produktivitas dalam dunia usaha dan industri.
\item Menghasilkan pengabdian kepada masyarakat melalui penerapan dan penyebarluasan ilmu pengetahuan dan teknologi secara profesional.
\item Terwujudnya kerja sama yang saling menguntungkan dengan berbagai pihak untuk meningkatkan daya saing.
\end{enumerate}

\section{Program Studi D4 Teknik Informatika}

\subsection{Visi}
Menjadi program studi yang unggul dalam bidang rekayasa perangkat lunak baik di tingkat nasional maupun internasional.

\subsection{Misi}
\begin{enumerate}
\item Melaksanakan pendidikan vokasi yang inovatif berdasarkan pada sistem pendidikan terapan dengan memanfaatkan kemajuan teknologi, sehingga mampu menghasilkan lulusan yang memiliki kompetensi di bidang rekayasa perangkat lunak dan siap bersaing di tingkat nasional dan global.
\item Melaksanakan penelitian terapan berbasis produk dan jasa bidang rekayasa perangkat lunak.
\item Melaksanakan pengabdian masyarakat dengan menggunakan kemajuan rekayasa perangkat lunak untuk meningkatkan kesejahteraan.
\item Mewujudkan kerja sama yang saling menguntungkan dengan berbagai pihak baik di dalam maupun di luar negeri pada bidang rekayasa perangkat lunak.
\end{enumerate}

\subsection{Tujuan}
\begin{enumerate}
\item Menghasilkan lulusan bidang rekayasa perangkat lunak yang berketuhanan, beretika dan bermoral baik, berpengetahuan dan berketrampilan tinggi, siap bekerja dan/atau berwirausaha yang mampu bersaing dalam skala nasional dan global.
\item Menghasilkan penelitian terapan bidang rekayasa perangkat lunak yang berskala nasional dan internasional.
\item Menghasilkan pengabdian kepada masyarakat melalui penerapan dan penyebarluasan ilmu pengetahuan dan teknologi secara profesional dalam bidang rekayasa perangkat lunak.
\item Terwujudnya kerja sama yang saling menguntungkan dengan berbagai pihak pada bidang rekayasa perangkat lunak untuk meningkatkan daya saing.
\end{enumerate}
